% from main.tex
%%%%%%%%%%  scan idx 108  |  folio 95  %%%%%%%%%%
Autrement dit, en abrégé, on peut prolonger un \(n\)-quasi-isomorphisme en un quasi-isomorphisme «~sans rien changer à ce qui est donné en degré \(\geq n\)~».

\section*{2. Complexes pseudo-cohérents}

\textbf{2.0.} Dans ce numéro, \(S\) désigne un site, \(S\) un objet de \(S\), \(\mathcal C\) une \(S\)-catégorie abélienne plate (1.1), \(\mathcal C_0\) une sous-\(S\)-catégorie additive strictement pleine de \(\mathcal C\). On suppose que \(\mathcal C_0\) est localement quasi-stable par noyau d'épimorphisme (1.2) et localement quasi-relevable dans \(\mathcal C\) (1.2).

\textbf{Définition 2.1.} \emph{Soient \(X\in\Ob S\), \(E\) un complexe de \(\mathcal C_X\), et \(n\in\mathbb Z\). Nous dirons que \(E\) est strictement \(\mathcal C_0\)-\(n\)-pseudo-cohérent (resp.\ \(\mathcal C_0\)-pseudo-cohérent, resp.\ \(\mathcal C_0\)-parfait) si \(E^i=0\) pour \(i\) assez grand et \(E^i\in\Ob\mathcal C_{0,S}\) pour \(i\geq n\) (resp.\ pour tout \(i\), resp.\ pour tout \(i\) et \(E^i=0\) pour presque tout \(i\)).}

Quand il n'y aura pas de confusion à craindre, nous dirons parfois «~strictement \(n\)-pseudo-cohérent (resp\dots)~» au lieu de «~strictement \(\mathcal C_0\)-\(n\)-pseudo-cohérent (resp\dots)~».

\textbf{Lemme 2.1.1.} \emph{Soient \(E\in\Ob\Cpx(\mathcal C_S)\) et \(n\in\mathbb Z\). Si \(E\) est strictement \(n\)-pseudo-cohérent et si \(\Ho^i(E)=0\) pour \(i>n\), alors \(\Ho^n(E)\) est de type fini.}

\emph{Preuve.} La suite
\[
  0\longrightarrow Z^n(E)\longrightarrow E^n\longrightarrow E^{n+1}\longrightarrow\cdots\longrightarrow 0
\]
est exacte et \(E^i\in\Ob\mathcal C_0\) pour \(i\geq n\). Comme \(\mathcal C_0\) est localement quasi-stable par noyau d'épimorphisme, il s'ensuit que \(Z^n(E)\) est de type fini, donc aussi \(\Ho^n(E)\), qui en est un quotient.

